\scanpage{166}
% archive foliation: 165
A K
\hfill\uline{Pseudomorphismes}

$X$ un préschéma

ouvert schématiquement dense $U$: si $X$ est le plus petit ss-préschéma fermé de $X$ \unsure{dominant|contenant} $U$ i.e. si $\overline U=X$ [lorsque $U\to X$ est quasi-compact cela implique que pour tout ouvert $V$, $U\cap V$ est schématiquement dense dans $X$ --- aussi est-ce \unsure{mieux comme} définition: \unsure{toute} \uline{\unsure{fonction|section}} de $\underline O_X$ sur un ouvert de $X$ nulle sur $U\cap V$ est nulle]
\margin{\unsure{définition équivalente}}

\deleted{\unsure{L'indépendance}}
Caractère local de la notion

$U$ schématiquement dense ds $X\Longleftrightarrow$ [si $X=\bigcup X_i$, les $X_i$ ouverts, les $U\cap X_i$ sch. denses dans les $X_i$]

$U$ schématiq. dense ds $X\Longrightarrow U$ schém. dense ds \unsure{des} $V$ q.c. $U\subset V\subset X$

Si $U,U'$ sont schém. denses, leur $U\cap U'$ l'est

[Soit $X$ et $Y$ sur $S$, \deleted{$\unsure{X|Y}$} $Y$ séparé sur $S$, \deleted{$X\Rightarrow$} $U$ schém. dense dans $X$. Si $f,g:X\to Y$ sont des $S$-morphismes tels que $f|U=g|U$

\deleted{\unsure{Quasi}}\inserted{Pseudo}-morphismes: classes d'équivalence d'\deleted{applications} morphismes définis sur des ouverts schématiquement denses, \unsure{aussi} $S$-morphismes. Pour un \unsure{pseudo}-morphisme dans un $S$-\unsure{pré}\uline{schéma}, il y a un domaine \inserted{(plus grand)} de définition.

\deleted{\unsure{le domaine} $\operatorname{Dom}(f|V)=$} $\operatorname{Dom}(f|V)=\operatorname{Dom}(f)\cap V$

